\documentclass[11pt,letterpaper]{article}
\usepackage[T1]{fontenc}
\usepackage[utf8]{inputenc}
\usepackage{lmodern}
\usepackage[margin=1in]{geometry}
\usepackage{amsmath,amssymb,amsthm}
\usepackage{microtype}
\usepackage[hidelinks]{hyperref}
\hypersetup{pdftitle={An elementary counterexample to Long's clique equality},
  pdfauthor={Alec Kriebel}}
\newcommand{\Z}{\mathbb Z}
\newcommand{\Q}{\mathbb Q}
\newcommand{\F}{\mathbb F}
\newcommand{\OO}{\mathcal O}
\newcommand{\Gu}{\Gamma_{\mathrm{unit}}}
\newcommand{\Gp}{\Gamma_{\mathrm{prime}}}
\newtheorem{theorem}{Theorem}
\newtheorem{lemma}{Lemma}

\title{An elementary counterexample to Long's clique equality}
\author{Alec Kriebel\\[2pt]
  \small Independent Researcher\\
  \small\href{https://orcid.org/0009-0001-9320-500X}
  {ORCID: 0009-0001-9320-500X}}
\date{17 October 2026}

\begin{document}
\maketitle

\begin{abstract}
For the full ring of integers of $\Q(\sqrt{10})$, both the unit-difference
graph and the graph allowing unit or prime-element differences have clique
number $2$, while the least norm of a nonzero prime ideal is $2$. Thus the
two sides of the equality asked about in Long's Question~7.2 in
arXiv:1512.04620v1 are $2$ and $4$. The proof uses the classical
nonprincipal ideal $(2,\sqrt{10})$ and its two-element residue field.
We include the elementary arithmetic and specify the prime-element
convention, which is essential to the conclusion. No novelty or historical
priority claim is made.
\end{abstract}

\section{The precise question}

Question~7.1 of \cite[p.~5]{problems} starts with an irreducible
$f\in\Z[X]$ all of whose roots are real, puts $k=\Q(\alpha)$ for a root
$\alpha$, and takes $\OO_k$ to be the full ring of integers. It asks how
often $\OO_k$ is a principal ideal domain; it does not impose that property
as a hypothesis. Question~7.2 \cite[p.~6]{problems} defines $\Gu$ on
$\OO_k$ by unit differences and enlarges it to $\Gp$ by allowing prime
differences. Its displayed formula is
\begin{equation}\label{eq:printed}
 |\operatorname{clique}(\Gp)|
 \leq |\operatorname{clique}(\Gu)|\,
       \min\{\,|\OO_k/I|\mid I\text{ prime}\,\},
\end{equation}
followed by the question ``Is this an equality?'' We write $\omega(G)$ for
the clique number denoted by $|\operatorname{clique}(G)|$ there.

Here a \emph{prime element} is a nonzero nonunit $p$ such that $(p)$ is a
prime ideal. This convention is fixed by the manuscript cited as
reference~[19] in the problem list: condition~(i) of its Definition~1
requires a difference to be a unit or to generate a prime ideal, and
Section~3 defines prime cliques using condition~(i) alone
\cite[pp.~1, 9]{long-ms}. The published Long--Thistlethwaite framework
retains these definitions \cite[Definition~1, p.~565; Section~3,
p.~572]{long}. The additional condition about surjectivity from the unit
group belongs to Lenstra--Hurwitz cliques and is not imposed on $\Gp$.

The minimum in \eqref{eq:printed} is over prime ideals. In contrast,
Theorem~3.1 of \cite[p.~573]{long} is printed with a minimum
\emph{principal}-ideal norm. We do not identify these minima or import
that theorem: the proof below evaluates each quantity in the question's
display directly.

\begin{theorem}\label{thm:main}
Let $k=\Q(\sqrt{10})$ and let $O=\OO_k$. With the definitions above,
\[
 \omega(\Gu)=\omega(\Gp)=2,
 \qquad
 \min_{\substack{\mathfrak q\text{ nonzero}\\
                  \text{prime ideal of }O}}|O/\mathfrak q|=2.
\]
Consequently the proposed equality in Long's Question~7.2 is false:
\[
 \omega(\Gp)=2<4
   =\omega(\Gu)
     \min_{\mathfrak q\text{ nonzero prime}}|O/\mathfrak q|.
\]
\end{theorem}

The ingredients are classical. Lenstra's residue argument already bounds
a unit clique by the cardinality of any proper ideal quotient
\cite[Proposition~2.2, p.~242]{lenstra}. Dummit's 2021 notes explicitly
give the ideal $(2,\sqrt{10})$, its nonprincipality, and the obstruction
to norm $\pm2$ modulo $5$ \cite[p.~13]{dummit}. These are credited
ingredients of the following self-contained calculation, rather than
claims of a new unit bound or new ideal arithmetic.

\section{The arithmetic and the two colors}

The polynomial $X^2-10$ is irreducible over $\Q$: a rational root of this
monic polynomial would be an integer with square $10$. Its two roots are
real, so this field satisfies the setup of Question~7.1.

\begin{lemma}\label{lem:ring}
The full ring of integers is $O=\Z[\sqrt{10}]$. For every nonzero
$\delta=a+b\sqrt{10}\in O$,
\begin{equation}\label{eq:index}
 [O:\delta O]=|a^2-10b^2|.
\end{equation}
\end{lemma}

\begin{proof}
Every $a+b\sqrt{10}$ with $a,b\in\Z$ is a root of the monic polynomial
$X^2-2aX+(a^2-10b^2)$ and is integral. Conversely, suppose
$r+s\sqrt{10}$ is integral, with $r,s\in\Q$. Its trace $2r$ and norm
$r^2-10s^2$ are integers. For $s\ne0$ this follows from the coefficients
of its monic minimal polynomial (by Gauss's lemma); for $s=0$ a rational
algebraic integer is an integer.

Put $m=2r\in\Z$. Since $40s^2=m^2-4(r^2-10s^2)\in\Z$, the reduced
denominator $v$ of $s$ satisfies $v^2\mid40$. Thus $v=1$ or $2$, and
$n=2s\in\Z$. Integrality of the norm now gives
\[
 m^2-10n^2\equiv0\pmod4.
\]
An odd $n$ would require a square to be $2$ modulo $4$. Hence $n$ is
even, and then $m$ is even too. Therefore $r,s\in\Z$.

In the integral basis $(1,\sqrt{10})$, multiplication by $\delta$ has
matrix
\[
 M_\delta=\begin{pmatrix}a&10b\\ b&a\end{pmatrix}.
\]
Its determinant is $a^2-10b^2\ne0$, since $\delta\ne0$ and
$\sqrt{10}\notin\Q$. Its image is $\delta O$. The index of the image
of a nonsingular integer matrix is the absolute value of its determinant:
unimodular row and column operations reduce to Smith normal form, where
this assertion is immediate. This proves \eqref{eq:index}.
\end{proof}

Define the surjective ring homomorphism
\[
 \varphi:O\longrightarrow\F_2,
 \qquad a+b\sqrt{10}\longmapsto a\pmod2.
\]
Multiplicativity follows because the constant coefficient of a product is
$ac+10bd$. The kernel is
\[
 P=(2,\sqrt{10}),
 \qquad O/P\cong\F_2.
\]
Indeed, a zero residue means that the constant coefficient is even.
Thus $P$ is a nonzero proper maximal prime ideal of norm $2$.

This ideal is nonprincipal. If $P=(a+b\sqrt{10})$, its index and
\eqref{eq:index} would give $a^2-10b^2=\pm2$. Modulo $5$ this requires
$a^2\equiv2$ or $3$, whereas the squares are $0,1,4$. This is impossible.

We next verify the maximality fact needed to exclude prime elements from
$P$. Let $Q$ be any nonzero prime ideal of $O$, and choose
$0\ne\delta\in Q$. By \eqref{eq:index}, $O/(\delta)$ is finite, so its
quotient $O/Q$ is a finite integral domain. Such a domain is a field:
multiplication by a nonzero element is injective, hence surjective, and
therefore has $1$ in its image. It follows that every nonzero prime ideal
of $O$ is maximal.

No unit lies in the proper ideal $P$, since otherwise $1\in P$.
If a prime element $p$ lay in $P$, the nonzero prime ideal $(p)$ would be
maximal and satisfy $(p)\subseteq P\subsetneq O$. Hence $(p)=P$, contrary
to nonprincipality. Thus $P$ contains neither units nor prime elements.

\begin{proof}[Proof of Theorem~\ref{thm:main}]
If two vertices have the same image under $\varphi$, their difference
lies in $P$ and cannot be an allowed edge difference in either graph.
The two residue classes therefore give a proper coloring of both graphs
with two colors. A clique has at most one vertex of each color, so both
clique numbers are at most $2$. The vertices $0$ and $1$ are adjacent in
both graphs, giving equality.

Every proper ideal quotient has at least two elements, since $0\ne1$ in
that quotient. The prime ideal $P$ attains this lower bound. Therefore
the minimum nonzero prime-ideal norm is $2$, and the two sides of the
proposed equality are $2$ and $2\cdot2=4$. If the unqualified word
``prime'' in \eqref{eq:printed} includes the zero ideal, its quotient is
the infinite ring $O$ and cannot lower this minimum.
\end{proof}

\section{Conventions, attribution, and verification}

The prime-element convention matters. The element $2$ is not prime:
$\sqrt{10}\notin(2)$ but $(\sqrt{10})^2=10\in(2)$. It is nevertheless
irreducible. Indeed, its absolute norm is $4$, absolute norm $1$ is
equivalent to being a unit (use conjugation for the inverse), and a
factorization into nonunits would require factors of absolute norm $2$,
which the congruence above excludes. Consequently the graph that allows
irreducible differences contains the triangle $\{0,1,2\}$ and is a
different problem.

Lopez's first prime-power graph $\mathcal P_k^{(1)}$ is exactly the
prime-element graph used here \cite[Definition~52, p.~26]{lopez}.
For $n\ge2$, the graph $\mathcal P_k^{(n)}$ also permits differences
whose principal ideals are higher powers of prime ideals. In this
example
\[
 P^2=(4,2\sqrt{10},10)=(2):
\]
all displayed generators are multiples of $2$, and $2=10-2\cdot4$.
Thus $\{0,1,2\}$ is a triangle in every $\mathcal P_k^{(n)}$ with
$n\ge2$, although $2$ is not a prime element. Lopez's bound $4$ in the
ramified case \cite[Theorem~55, p.~27]{lopez} concerns
$\mathcal P_k^{(1)}$ and is consistent with its exact value $2$ here.
The ramification clause just cited is the thesis's wording; the journal
Corollary~3.1.1 \cite[p.~573]{long} literally states its $4$-element
bound for the case where $2$ splits.

The conclusion concerns the equality exactly printed in
\cite[Question~7.2]{problems}, with the full ring of integers and no
principal-ideal-domain premise. It makes no assertion about a separately
restricted PID version. The preceding citations record the definitions
and classical ingredients used here. We make no claim of novelty,
firstness, historical priority, or current worldwide openness of the
printed question.

An accompanying verification package records source receipts and exact
finite diagnostic checks of the elementary arithmetic and boundary
conventions. Those computations are supplementary; the proof above is
complete and does not infer universal claims from a finite search.
Third-party primary-source PDFs are cited rather than redistributed.

\paragraph{Disclosure.}
AI tools were used extensively in solving, verifying, and writing this
note. The note is unrefereed, has received no external human peer review,
and has no formal proof-assistant certification.

\begin{thebibliography}{9}

\bibitem{problems}
K.~Delp, D.~Hoffoss, and J.~F.~Manning (eds.),
\emph{Problems In Groups, Geometry, and Three-Manifolds},
arXiv:1512.04620v1 (2015), Questions~7.1--7.2, pp.~5--6;
reference~[19], p.~7.
\url{https://arxiv.org/abs/1512.04620v1}.

\bibitem{long-ms}
D.~D.~Long and M.~B.~Thistlethwaite,
\emph{Lenstra-Hurwitz cliques and the class number one problem},
16-page author manuscript cited in \cite{problems},
Definition~1, p.~1; Section~3, pp.~9--10.
\url{https://web.math.ucsb.edu/~long/pubpdf/LenstraHurwitz.pdf}.

\bibitem{long}
D.~D.~Long and M.~B.~Thistlethwaite,
Lenstra--Hurwitz cliques and the class number one problem,
\emph{Journal of Number Theory} \textbf{162} (2016), 564--577.
\href{https://doi.org/10.1016/j.jnt.2015.11.002}
{doi:10.1016/j.jnt.2015.11.002}.

\bibitem{lenstra}
H.~W.~Lenstra, Jr.,
Euclidean number fields of large degree,
\emph{Inventiones Mathematicae} \textbf{38} (1977), 237--254,
especially Proposition~2.2, p.~242.
\url{https://pub.math.leidenuniv.nl/~lenstrahw/PUBLICATIONS/1977b/art.pdf}.

\bibitem{dummit}
E.~Dummit,
\emph{Number Theory (part 9): The Geometry of Numbers},
2021, version~1.90, p.~13 (the $\Z[\sqrt{10}]$ example).
\url{https://dummit.cos.northeastern.edu/teaching_sp21_4527/numthy_9_geometry_of_numbers_v1.90.pdf}.

\bibitem{lopez}
D.~S.~Lopez,
\emph{Lenstra-Hurwitz Cliques In Real Quadratic Fields},
M.S. thesis, San Jos\'e State University, 2018,
Definition~52, p.~26; Theorem~55, p.~27.
\href{https://doi.org/10.31979/etd.3bc5-77v2}
{doi:10.31979/etd.3bc5-77v2}.

\end{thebibliography}
\end{document}
